\documentclass[11pt]{article}
\usepackage[margin=1in]{geometry}
\usepackage{amsmath,amssymb}
\usepackage[hidelinks]{hyperref}

\newcommand{\R}{\mathbb R}
\newcommand{\K}{\mathbb K}

\begin{document}

\begin{center}
{\Large Literature-implied answer for arXiv:0711.3768}\\[4pt]
{\normalsize Removing the infinite-dimensional Hilbert restriction in the
convex-transitive \(C_0(L,H)\) theorem}
\end{center}

\section*{Status}

This is a lightweight literature-implied-answer packet.  It records that a
later theorem answers a source-paper question after direct identification.  It
is not a new proof packet.

\section*{Original Question}

In Section 2 of Talponen's \emph{Convex-transitivity and function spaces}
\cite{Talponen2007}, in the subsection on \(C_0(L,X)\), the paper proves the
following two-part theorem:
\begin{enumerate}
\item if \(X\) is strictly convex and \(C_0(L,X)\) is convex-transitive, then
  \(X\) is convex-transitive;
\item if \(H\) is a Hilbert space, \(\dim(H)=\infty\), and scalar \(C_0(L)\)
  is convex-transitive, then \(C_0(L,H)\) is convex-transitive.
\end{enumerate}
Immediately after the theorem, the source says that the proof relies strongly
on \(\dim(H)=\infty\), and that it is not known whether this assumption can be
removed.

\section*{Supporting Theorem}

Section 4, Theorem 4.3 of Rambla-Barreno and Talponen's
\emph{Uniformly convex-transitive function spaces} \cite{RT2009} states:
if \(C_0^{\R}(L)\) is convex-transitive and \(X\) is a convex-transitive Banach
space over \(\K\), then \(C_0^{\K}(L,X)\) is convex-transitive.

\section*{Identification}

Take \(X=H\).  Every Hilbert space is convex-transitive: for any two unit
vectors, a surjective linear isometry sends one to the other, so every unit
vector is a big point.  Thus Theorem 4.3 of \cite{RT2009} applies to every
Hilbert space \(H\), including finite-dimensional \(H\).  The case
\(\dim(H)=1\) is the scalar \(C_0(L)\) case itself.

Therefore the infinite-dimensional Hilbert hypothesis in the source theorem is
not needed.  The relation is literature-implied rather than explicitly
announced: the 2009 theorem is broader and does not appear to single out the
exact 2007 sentence.

\section*{Scope}

This packet answers only the source-paper question about removing
\(\dim(H)=\infty\) from the Hilbert-valued \(C_0(L,H)\) convex-transitivity
theorem.  It does not address unrelated open questions in the same source,
such as the separate question about \(L^\infty(X)\) and convex-transitivity.

\begin{thebibliography}{9}
\raggedright
\bibitem{Talponen2007}
J. Talponen,
\emph{Convex-transitivity and function spaces},
arXiv:0711.3768v2.

\bibitem{RT2009}
F. Rambla-Barreno and J. Talponen,
\emph{Uniformly convex-transitive function spaces},
arXiv:0902.0640v2.
\end{thebibliography}

\end{document}
